% Lapisan solusi O001 -- Bab 8 (Spektrum)
% 2 solusi asli terpisah; bukan tulisan John M. Erdman.
% Provenans model: OpenAI Codex gpt-5.6-sol, Ultra, atas arahan pengguna.
% Lisensi komponen solusi: CC BY-SA 4.0.
% Tidak menyiratkan dukungan John M. Erdman atau Portland State University.

\markboth{SOLUSI O001 UNTUK BAB 8: SPEKTRUM}{SOLUSI O001 UNTUK BAB 8: SPEKTRUM}
\section*{Solusi O001 untuk Bab 8: Spektrum}
\addcontentsline{toc}{section}{Solusi O001 untuk Bab 8: Spektrum}
\begin{center}
\small
Komponen solusi asli terpisah, CC BY-SA 4.0. Ditulis dengan bantuan
OpenAI Codex gpt-5.6-sol, Ultra, atas arahan pengguna. Pernyataan latihan
berasal dari edisi terjemahan karya John M. Erdman; solusi ini bukan tulisan
beliau dan tidak menyiratkan dukungan beliau atau Portland State University.
\end{center}

% O001-SOLUTION-ID: O001-FAOA-2015-CH08-EX-001-SOLUTION
% SOURCE-EXERCISE-ID: FAOA-2015-CH08-NODE-0031
% STATEMENT-TARGET-SHA256: e16770f5d0da46b3543a0b001816b53299c752a00a3a6db2d47672ff3d0309cf
\begin{o001solution}
  {O001-FAOA-2015-CH08-EX-001-SOLUTION}
  {FAOA-2015-CH08-NODE-0031}
  {e16770f5d0da46b3543a0b001816b53299c752a00a3a6db2d47672ff3d0309cf}
\begin{o001statement}
Identifikasi hasil kali dan koproduk (jika ada) dalam kategori $\cat{BALG}$
yang objeknya aljabar Banach
 \index{hasil kali!dalam $\cat{BALG}$}%
 \index{Banach!aljabar!hasil kali}%
 \index{Banach!aljabar!koproduk}%
 \index{koproduk!dalam $\cat{BALG}$}%
 \index{BALG@$\cat{BALG}$!hasil kali dan koproduk dalam}%
dan morfismenya homomorfisme aljabar kontinu, serta dalam kategori yang objeknya aljabar Banach dan
morfismenya homomorfisme aljabar kontraktif. (Di sini,
 \index{kontraktif}%
\emph{kontraktif} berarti $\norm{Tx} \le \norm x$ untuk setiap~$x$.)
\end{o001statement}
\begin{o001answer}
Dalam $\cat{BALG}$ dengan homomorfisme kontinu, hasil kali $A$ dan $B$
adalah $A\times B$ dengan operasi titik demi titik dan, misalnya, norma
$\norm{(a,b)}=\norm a+\norm b$; kategori ini tidak mempunyai koproduk
biner pada umumnya. Dalam kategori homomorfisme kontraktif, hasil kali
adalah $A\times B$ dengan norma maksimum, sedangkan koproduknya adalah
produk bebas Banach kontraktif $A*B$, yaitu penyelesaian jumlah-$l_1$
dari kata-kata tensor berselang-seling dalam $A$ dan~$B$.
\end{o001answer}
\begin{o001proof}
Untuk homomorfisme kontinu, proyeksi koordinat dari $A\times B$ kontinu.
Jika $f:C\to A$ dan $g:C\to B$ kontinu, satu-satunya pemetaan yang membuat
diagram hasil kali komutatif ialah
\[
  h:C\longrightarrow A\times B,\qquad h(c)=(f(c),g(c)).
\]
Pemetaan ini homomorfisme dan
$\norm{h(c)}\le(\norm f+\norm g)\norm c$, sehingga kontinu. Ini membuktikan
sifat universal hasil kali. Norma maksimum juga memberi objek hasil kali
yang isomorfik dalam kategori ini.

Untuk melihat bahwa koproduk tidak selalu ada dalam kategori morfisme
kontinu, ambil aljabar Banach satu dimensi $E=\C$ dengan perkalian nol.
Andaikan $(P,i,j)$ adalah koproduk dua salinan $E$, dan tulis
$p=i(1)$ serta $q=j(1)$. Untuk setiap $\lambda>0$, pemetaan
\[
 f_\lambda(z)=\lambda z e_{12},\qquad
 g_\lambda(z)=\lambda z e_{21}
\]
adalah homomorfisme kontinu $E\to M_2(\C)$ karena
$e_{12}^2=e_{21}^2=0$. Sifat universal akan memberi homomorfisme kontinu
$H_\lambda:P\to M_2(\C)$ dengan nilai-nilai tersebut. Maka
\[
 H_\lambda((pq)^n)=\lambda^{2n}e_{11}.
\]
Dengan norma operator pada $M_2(\C)$ dan submultiplikativitas norma di
$P$, kontinuitas $H_\lambda$ akan memaksa
\[
 \lambda^{2n}\le\norm{H_\lambda}
       (\norm p\,\norm q)^n\qquad(n\ge1).
\]
Pemetaan universal ke aljabar perkalian-nol menunjukkan $p,q\ne0$.
Memilih $\lambda^2>\norm p\,\norm q$ membuat ketaksamaan terakhir mustahil
saat $n\to\infty$. Jadi kategori ini tidak mempunyai semua koproduk biner.

Dalam kategori homomorfisme kontraktif, beri $A\times B$ norma
\[
  \norm{(a,b)}_{\max}=\max\{\norm a,\norm b\}.
\]
Proyeksinya kontraktif, dan pasangan $h=(f,g)$ kontraktif apabila $f$ dan
$g$ kontraktif. Jadi inilah hasil kali pada kategori tersebut.

Untuk koproduk kontraktif, bagi setiap kata takkosong berselang-seling
$\epsilon=(\epsilon_1,\ldots,\epsilon_n)$ dengan huruf $A,B$, bentuk hasil
kali tensor projektif
\[
 X_{\epsilon_1}\mathbin{\widehat{\otimes}_{\pi}}\cdots
 \mathbin{\widehat{\otimes}_{\pi}}X_{\epsilon_n},
 \qquad X_A=A,\quad X_B=B,
\]
lalu ambil jumlah langsung-$l_1$ atas semua kata tersebut. Perkalian
dilakukan dengan menggabungkan kata; jika dua huruf pada sambungan sama,
kedua faktor batas dikalikan di dalam $A$ atau $B$. Submultiplikativitas
norma aljabar dan sifat norma projektif membuat operasi ini kontraktif dan
memanjang ke penyelesaiannya, yang kita notasikan $A*B$.

Penyertaan ke suku satu-huruf bersifat isometrik. Jika
$f:A\to C$ dan $g:B\to C$ kontraktif, pada tensor elementer definisikan
\[
 x_1\otimes\cdots\otimes x_n\longmapsto
 F_{\epsilon_1}(x_1)\cdots F_{\epsilon_n}(x_n),
 \qquad F_A=f,\quad F_B=g.
\]
Norma hasilnya paling besar daripada hasil kali norma faktor, sehingga
pemetaan pada setiap suku dan pada jumlah-$l_1$ kontraktif. Pemetaan ini
adalah homomorfisme dan memanjang secara unik ke $A*B$; keunikannya berlaku
karena kata-kata dalam kedua citra penyertaan rapat. Jadi $A*B$ memenuhi
sifat universal koproduk kontraktif.
\end{o001proof}
\end{o001solution}

% O001-SOLUTION-ID: O001-FAOA-2015-CH08-EX-002-SOLUTION
% SOURCE-EXERCISE-ID: FAOA-2015-CH08-NODE-0064
% STATEMENT-TARGET-SHA256: 8542ce320f7fb680060f8d11f114b381d061752b828a1f4a5e5f5080a8287e9c
\begin{o001solution}
  {O001-FAOA-2015-CH08-EX-002-SOLUTION}
  {FAOA-2015-CH08-NODE-0064}
  {8542ce320f7fb680060f8d11f114b381d061752b828a1f4a5e5f5080a8287e9c}
\begin{o001statement}
Dengan menggunakan rumus untuk operator Volterra $V$ dalam contoh~\ref{000319} sebagai
 \index{operator Volterra!dan persamaan integral}%
 \index{operator!Volterra!dan persamaan integral}%
acuan, pada ruang Banach~$\fml C([0,1])$ definisikan operator integral melalui rumus $Vf(x)=\int_0^x f(t)\,dt$.
 \begin{enumerate}
  \item[(a)] Hitung spektrum operator $V$. (\emph{Petunjuk.} Tunjukkan bahwa $\norm{V^n}
\le (n!)^{-1}$ untuk setiap $n \in \N$ dan gunakan \emph{rumus radius spektral}.)
  \item[(b)] Apa implikasi (a) terhadap kemungkinan memperoleh fungsi kontinu $f$ yang
memenuhi persamaan integral
   \[ f(x) - \mu\int_0^xf(t)\,dt = h(x), \]
dengan $\mu$ suatu skalar dan $h$ suatu fungsi kontinu yang diberikan pada~$[0,1]$?
  \item[(c)] Gunakan gagasan dalam (b) dan uraian deret Neumann untuk $\vc 1 - \mu V$ (lihat
proposisi~\ref{prop_Neumann_series}) untuk menghitung secara eksplisit (dan dalam bentuk fungsi elementer)
suatu solusi untuk persamaan integral
 \[f(x) - {\textstyle \frac12}\int_0^xf(t)\,dt = e^x.\]
 \end{enumerate}
\end{o001statement}
\begin{o001answer}
\begin{enumerate}
 \item[(a)] $\sigma(V)=\{0\}$.
 \item[(b)] Untuk setiap skalar $\mu$ dan setiap
 $h\in\fml C([0,1])$, terdapat tepat satu solusi, yaitu
 \[
   f=(\vc1-\mu V)^{-1}h=\sum_{n=0}^{\infty}\mu^nV^nh.
 \]
 \item[(c)] Solusi eksplisitnya adalah
 $f(x)=2e^x-e^{x/2}$.
\end{enumerate}
\end{o001answer}
\begin{o001proof}
Induksi dan teorema Fubini untuk integral kontinu memberi, untuk $n\ge1$,
\[
 (V^nf)(x)=\frac1{(n-1)!}\int_0^x(x-t)^{n-1}f(t)\,dt.
\]
Karena $0\le x\le1$,
\[
 |(V^nf)(x)|\le\frac{x^n}{n!}\norm f_u
 \le\frac1{n!}\norm f_u,
\]
sehingga $\norm{V^n}\le1/n!$. Rumus radius spektral menghasilkan
\[
 \rho(V)\le\lim_{n\to\infty}(1/n!)^{1/n}=0.
\]
Spektrum operator pada ruang Banach kompleks takkosong, jadi
$\sigma(V)\subseteq\{0\}$ memaksa $\sigma(V)=\{0\}$. Secara langsung,
$0$ memang berada dalam spektrum karena $V$ tidak surjektif: setiap fungsi
dalam jangkauannya bernilai nol di~$0$.

Untuk setiap $\mu$, $\sigma(\mu V)=\{0\}$ (termasuk $\mu=0$), sehingga
$\vc1-\mu V$ invertibel. Lebih konkret, estimasi di atas memberi
\[
 \sum_{n=0}^{\infty}\norm{\mu^nV^n}
 \le\sum_{n=0}^{\infty}\frac{|\mu|^n}{n!}=e^{|\mu|}<\infty.
\]
Karena itu deret operator $R=\sum_{n\ge0}\mu^nV^n$ konvergen dalam norma.
Mengalikan jumlah parsial dengan $\vc1-\mu V$ menghasilkan
$\vc1-\mu^{N+1}V^{N+1}$, yang menuju $\vc1$ dalam norma; urutan perkalian
yang lain sama. Maka $R=(\vc1-\mu V)^{-1}$ dan untuk setiap $h$ terdapat
tepat satu solusi $f=Rh$.

Untuk $h(x)=e^x$ dan $n\ge1$, induksi dari integrasi menunjukkan
\[
 (V^ne^\cdot)(x)=e^x-\sum_{k=0}^{n-1}\frac{x^k}{k!}.
\]
Dengan $\mu=1/2$, penataan ulang sah karena konvergensi mutlak dan seragam,
dan
\begin{align*}
 f(x)
 &=e^x+\sum_{n=1}^{\infty}2^{-n}
       \left(e^x-\sum_{k=0}^{n-1}\frac{x^k}{k!}\right)\\
 &=2e^x-\sum_{k=0}^{\infty}\frac{x^k}{k!}
       \sum_{n=k+1}^{\infty}2^{-n}
 =2e^x-\sum_{k=0}^{\infty}\frac{(x/2)^k}{k!}
 =2e^x-e^{x/2}.
\end{align*}
Sebagai pemeriksaan langsung,
\[
 \frac12\int_0^x(2e^t-e^{t/2})\,dt=e^x-e^{x/2},
\]
sehingga $f(x)-\frac12\int_0^xf(t)\,dt=e^x$.
\end{o001proof}
\end{o001solution}
